% 1 John — Generated by generate_bible.py
% NET Bible text, scripture package formatting
\gdef\bbookchaptable{\vspace{4pt}%
\parbox[t]{\linewidth}{\centering\small\hyperlink{ch-1john-1}{\textit{\small 1}}\hspace{0.4em}\hyperlink{ch-1john-2}{\textit{\small 2}}\hspace{0.4em}\hyperlink{ch-1john-3}{\textit{\small 3}}\hspace{0.4em}\hyperlink{ch-1john-4}{\textit{\small 4}}\hspace{0.4em}\hyperlink{ch-1john-5}{\textit{\small 5}}}%
\vspace{6pt}}
\bbook{The First Epistle General of John}{}{}{1john}

\begin{scripture}

\markboth{1 John 1:1}{1 John 1:1}
\ch{1} \allowchapbreak\hypertarget{ch-1john-1}{}\lettrine[lines=8]{\color{general}T}{\textsc{his}} is what we proclaim to you: what was from the beginning, what we have heard, what we have seen with our eyes, what we have looked at and our hands have touched (concerning the word of life---\gva{a}\marginnote{\gva{a}\,I heard him speak, I saw him myself with my eyes, I handled with my hands him that is true God, being made true man, and not I alone, but others also that were with me.}\gva{b}\marginnote{\gva{b}\,That same everlasting Word by whom all things are made, and in whom only is there life.}\everypar{}
\markboth{1 John 1:2}{1 John 1:2}\vs{2} and the life was revealed, and we have seen and testify and announce to you the eternal life that was with the Father and was revealed to us).\gva{c}\marginnote{\gva{c}\,Being sent by him: and that doctrine is correctly said to be shown, for no man could so much as have thought of it, if it had not been thus shown.}
\markboth{1 John 1:3}{1 John 1:3}\vs{3} What we have seen and heard we announce to you too, so that you may have fellowship with us (and indeed our fellowship is with the Father and with his Son Jesus Christ).
\markboth{1 John 1:4}{1 John 1:4}\vs{4} Thus we are writing these things so that our joy may be complete.
\everypar{}

\parshape=0
\markboth{1 John 1:5}{1 John 1:5}\vs{5} Now this is the gospel message we have heard from him and announce to you: God is light, and in him there is no darkness at all.
\markboth{1 John 1:6}{1 John 1:6}\vs{6} If we say we have fellowship with him and yet keep on walking in the darkness, we are lying and not practicing the truth.
\markboth{1 John 1:7}{1 John 1:7}\vs{7} But if we walk in the light as he himself is in the light, we have fellowship with one another and the blood of Jesus his Son cleanses us from all sin.\gva{d}\marginnote{\gva{d}\,God is said to be light by his own nature, and to be in light, that is to say, in that everlasting infinite blessedness: and we are said to walk in light in that the beams of that light shine to us in the Word.}
\markboth{1 John 1:8}{1 John 1:8}\vs{8} If we say we do not bear the guilt of sin, we are deceiving ourselves and the truth is not in us.\gva{e}\marginnote{\gva{e}\,This fully refutes that perfectness of works of supererogation (doing more than duty requires, the idea that excess good works can form a reserve fund of merit that can be drawn on in favour of sinners) which the papists dream of.}\gva{f}\marginnote{\gva{f}\,So then, John speaks not thus for modesty's sake, as some say but because it is so indeed.}
\markboth{1 John 1:9}{1 John 1:9}\vs{9} But if we confess our sins, he is faithful and righteous, forgiving us our sins and cleansing us from all unrighteousness.\gva{g}\marginnote{\gva{g}\,So then our salvation depends on the free promise of God, who because he is faithful and just, will perform that which he hath promised.}\gva{h}\marginnote{\gva{h}\,Where then are our merits? for this is our true happiness.}
\markboth{1 John 1:10}{1 John 1:10}\vs{10} If we say we have not sinned, we make him a liar and his word is not in us.\gva{i}\marginnote{\gva{i}\,They do not only deceive themselves but are blasphemous against God.}\gva{k}\marginnote{\gva{k}\,His doctrine shall have no place in us; that is, in our hearts.}

\markboth{1 John 2:1}{1 John 2:1}
\Needspace*{8\baselineskip}
\ch{2} \allowchapbreak\hypertarget{ch-1john-2}{}(\lettrine[lines=5]{\color{general}M}{\textsc{y}} little children, I am writing these things to you so that you may not sin.) But if anyone does sin, we have an advocate with the Father, Jesus Christ the Righteous One,\gva{a}\marginnote{\gva{a}\,In that be names Christ, he eliminates all others.}\everypar{}
\markboth{1 John 2:2}{1 John 2:2}\vs{2} and he himself is the atoning sacrifice for our sins, and not only for our sins but also for the whole world.\gva{b}\marginnote{\gva{b}\,Reconciliation and intercession go together, to give us to understand that he is both advocate and high priest.}\gva{c}\marginnote{\gva{c}\,For men of all sorts, of all ages, and all places, so that this benefit being not to the Jews only, of whom he speaks as appears in 1 John 2:7 but also to other nations.}
\\\indent\markboth{1 John 2:3}{1 John 2:3}\vs{3} Now by this we know that we have come to know God: if we keep his commandments.\gva{d}\marginnote{\gva{d}\,This must be understood of such a knowledge as has faith with it, and not of a common knowledge.}\gva{e}\marginnote{\gva{e}\,For the tree is known by the fruit.}
\markboth{1 John 2:4}{1 John 2:4}\vs{4} The one who says ``I have come to know God'' and yet does not keep his commandments is a liar, and the truth is not in such a person.
\markboth{1 John 2:5}{1 John 2:5}\vs{5} But whoever obeys his word, truly in this person the love of God has been perfected. By this we know that we are in him.\gva{f}\marginnote{\gva{f}\,Wherewith we love God.}\gva{g}\marginnote{\gva{g}\,He means our union with Christ.}
\markboth{1 John 2:6}{1 John 2:6}\vs{6} The one who says he resides in God ought himself to walk just as Jesus walked.
\everypar{}

\parshape=0
\markboth{1 John 2:7}{1 John 2:7}\vs{7} Dear friends, I am not writing a new commandment to you, but an old commandment which you have had from the beginning. The old commandment is the word that you have already heard.
\markboth{1 John 2:8}{1 John 2:8}\vs{8} On the other hand, I am writing a new commandment to you, which is true in him and in you, because the darkness is passing away and the true light is already shining.
\markboth{1 John 2:9}{1 John 2:9}\vs{9} The one who says he is in the light but still hates his fellow Christian is still in the darkness.
\markboth{1 John 2:10}{1 John 2:10}\vs{10} The one who loves his fellow Christian resides in the light, and there is no cause for stumbling in him.
\markboth{1 John 2:11}{1 John 2:11}\vs{11} But the one who hates his fellow Christian is in the darkness, walks in the darkness, and does not know where he is going because the darkness has blinded his eyes.
\everypar{}

\parshape=0
\markboth{1 John 2:12}{1 John 2:12}\vs{12} I am writing to you, little children, that your sins have been forgiven because of his name.\gva{i}\marginnote{\gva{i}\,Therefore I write to you, because you are of their number whom God has reconciled to himself.}\gva{k}\marginnote{\gva{k}\,For his own sake: in that he names Christ he eliminates all others, whether they are in heaven or on earth.}
\markboth{1 John 2:13}{1 John 2:13}\vs{13} I am writing to you, fathers, that you have known him who has been from the beginning. I am writing to you, young people, that you have conquered the evil one.
\markboth{1 John 2:14}{1 John 2:14}\vs{14} I have written to you, children, that you have known the Father. I have written to you, fathers, that you have known him who has been from the beginning. I have written to you, young people, that you are strong, and the word of God resides in you, and you have conquered the evil one.
\everypar{}

\parshape=0
\markboth{1 John 2:15}{1 John 2:15}\vs{15} Do not love the world or the things in the world. If anyone loves the world, the love of the Father is not in him,\gva{l}\marginnote{\gva{l}\,He speaks of the world, as it agrees not with the will of God, for otherwise God is said to love the world with an infinite love, John 3:16 that is to say, those whom he chose out of the world.}\gva{m}\marginnote{\gva{m}\,Wherewith the Father is loved.}
\markboth{1 John 2:16}{1 John 2:16}\vs{16} because all that is in the world (the desire of the flesh and the desire of the eyes and the arrogance produced by material possessions) is not from the Father, but is from the world.
\markboth{1 John 2:17}{1 John 2:17}\vs{17} And the world is passing away with all its desires, but the person who does the will of God remains forever.
\everypar{}

\parshape=0
\markboth{1 John 2:18}{1 John 2:18}\vs{18} Children, it is the last hour, and just as you heard that the antichrist is coming, so now many antichrists have appeared. We know from this that it is the last hour.\gva{n}\marginnote{\gva{n}\,He uses this word "Little" not because he speaks to children, but to allure them the more by using such sweet words.}
\markboth{1 John 2:19}{1 John 2:19}\vs{19} They went out from us, but they did not really belong to us because if they had belonged to us, they would have remained with us. But they went out from us to demonstrate that all of them do not belong to us.\gva{o}\marginnote{\gva{o}\,So then the elect can never fall from grace.}
\everypar{}

\parshape=0
\markboth{1 John 2:20}{1 John 2:20}\vs{20} Nevertheless you have an anointing from the Holy One, and you all know.\gva{p}\marginnote{\gva{p}\,The grace of the Holy Spirit, and this is a borrowed type of speech taken from the anointings used in the law.}\gva{q}\marginnote{\gva{q}\,From Christ who is peculiarly called Holy.}
\markboth{1 John 2:21}{1 John 2:21}\vs{21} I have not written to you that you do not know the truth, but that you do know it, and that no lie is of the truth.
\markboth{1 John 2:22}{1 John 2:22}\vs{22} Who is the liar but the person who denies that Jesus is the Christ? This one is the antichrist: the person who denies the Father and the Son.\gva{r}\marginnote{\gva{r}\,Is the true Messiah.}
\markboth{1 John 2:23}{1 John 2:23}\vs{23} Everyone who denies the Son does not have the Father either. The person who confesses the Son has the Father also.\gva{s}\marginnote{\gva{s}\,They deceive themselves, and also deceive others, who say that the Moslems and other infidels worship the same God that we do.}
\everypar{}

\parshape=0
\markboth{1 John 2:24}{1 John 2:24}\vs{24} As for you, what you have heard from the beginning must remain in you. If what you heard from the beginning remains in you, you also will remain in the Son and in the Father.
\markboth{1 John 2:25}{1 John 2:25}\vs{25} Now this is the promise that he himself made to us: eternal life.
\markboth{1 John 2:26}{1 John 2:26}\vs{26} These things I have written to you about those who are trying to deceive you.
\everypar{}

\parshape=0
\markboth{1 John 2:27}{1 John 2:27}\vs{27} Now as for you, the anointing that you received from him resides in you, and you have no need for anyone to teach you. But as his anointing teaches you about all things, it is true and is not a lie. Just as it has taught you, you reside in him.\gva{t}\marginnote{\gva{t}\,The Spirit who you have received from Christ, and who has led you into all truth.}\gva{u}\marginnote{\gva{u}\,You are not ignorant of those things, and therefore I teach them not as things that were never heard of, but call them to your mind as things which you do know.}\gva{x}\marginnote{\gva{x}\,He commends both the doctrine which they had embraced, and also highly praises their faith, and the diligence of those who taught them, yet so, that he takes nothing from the honour due to the Holy Spirit.}
\everypar{}

\parshape=0
\markboth{1 John 2:28}{1 John 2:28}\vs{28} And now, little children, remain in him, so that when he appears we may have confidence and not shrink away from him in shame when he comes back.
\markboth{1 John 2:29}{1 John 2:29}\vs{29} If you know that he is righteous, you also know that everyone who practices righteousness has been fathered by him.

\markboth{1 John 3:1}{1 John 3:1}
\Needspace*{8\baselineskip}
\ch{3} \allowchapbreak\hypertarget{ch-1john-3}{}(\lettrine[lines=5]{\color{general}S}{\textsc{ee}} what sort of love the Father has given to us: that we should be called God's children---and indeed we are! For this reason the world does not know us: Because it did not know him.\gva{a}\marginnote{\gva{a}\,What a gift of how great love.}\gva{b}\marginnote{\gva{b}\,That we should be the sons of God, and so, that all the world may see that we are so.}\everypar{}
\markboth{1 John 3:2}{1 John 3:2}\vs{2} Dear friends, we are God's children now, and what we will be has not yet been revealed. We know that whenever it is revealed we will be like him because we will see him just as he is.\gva{c}\marginnote{\gva{c}\,Like, but not equal.}\gva{d}\marginnote{\gva{d}\,For now we see as in a glass 1 Corinthians 13:12}
\markboth{1 John 3:3}{1 John 3:3}\vs{3} And everyone who has this hope focused on him purifies himself, just as Jesus is pure).\gva{e}\marginnote{\gva{e}\,This word signifies a likeness, but not an equality.}
\everypar{}

\parshape=0
\markboth{1 John 3:4}{1 John 3:4}\vs{4} Everyone who practices sin also practices lawlessness; indeed, sin is lawlessness.\gva{f}\marginnote{\gva{f}\,Does not give himself to pureness.}\gva{g}\marginnote{\gva{g}\,A short definition of sin.}
\markboth{1 John 3:5}{1 John 3:5}\vs{5} And you know that Jesus was revealed to take away sins, and in him there is no sin.
\markboth{1 John 3:6}{1 John 3:6}\vs{6} Everyone who resides in him does not sin; everyone who sins has neither seen him nor known him.\gva{h}\marginnote{\gva{h}\,He is said to sin, that does not give himself to purity, and in him sin reigns: but sin is said to dwell in the faithful, and not to reign in them.}
\markboth{1 John 3:7}{1 John 3:7}\vs{7} Little children, let no one deceive you: The one who practices righteousness is righteous, just as Jesus is righteous.
\markboth{1 John 3:8}{1 John 3:8}\vs{8} The one who practices sin is of the devil because the devil has been sinning from the beginning. For this purpose the Son of God was revealed: to destroy the works of the devil.\gva{i}\marginnote{\gva{i}\,Resembles the devil, as the child does the father, and is governed by his Spirit.}\gva{k}\marginnote{\gva{k}\,He says not "sinned" but "sins" for he does nothing else but sin.}\gva{l}\marginnote{\gva{l}\,From the very beginning of the world.}
\markboth{1 John 3:9}{1 John 3:9}\vs{9} Everyone who has been fathered by God does not practice sin because God's seed resides in them, and thus they are not able to sin because they have been fathered by God.\gva{m}\marginnote{\gva{m}\,The Holy Spirit is so called by the effect he works, because by his power and mighty working, as it were by seed, we are made new men.}
\markboth{1 John 3:10}{1 John 3:10}\vs{10} By this the children of God and the children of the devil are revealed: Everyone who does not practice righteousness---the one who does not love his fellow Christian---is not of God.
\everypar{}

\parshape=0
\markboth{1 John 3:11}{1 John 3:11}\vs{11} For this is the gospel message that you have heard from the beginning: that we should love one another,
\markboth{1 John 3:12}{1 John 3:12}\vs{12} not like Cain who was of the evil one and brutally murdered his brother. And why did he murder him? Because his deeds were evil, but his brother's were righteous.
\everypar{}

\parshape=0
\markboth{1 John 3:13}{1 John 3:13}\vs{13} Therefore do not be surprised, brothers and sisters, if the world hates you.
\markboth{1 John 3:14}{1 John 3:14}\vs{14} We know that we have crossed over from death to life because we love our fellow Christians. The one who does not love remains in death.\gva{o}\marginnote{\gva{o}\,Love is a token that we are translated from death to life, for by the effects the cause is known.}
\markboth{1 John 3:15}{1 John 3:15}\vs{15} Everyone who hates his fellow Christian is a murderer, and you know that no murderer has eternal life residing in him.
\markboth{1 John 3:16}{1 John 3:16}\vs{16} We have come to know love by this: that Jesus laid down his life for us; thus we ought to lay down our lives for our fellow Christians.
\markboth{1 John 3:17}{1 John 3:17}\vs{17} But whoever has the world's possessions and sees his fellow Christian in need and shuts off his compassion against him, how can the love of God reside in such a person?\gva{p}\marginnote{\gva{p}\,Wherewith this life is sustained.}\gva{q}\marginnote{\gva{q}\,Opens not his heart to him, nor helps him willingly and cheerfully.}
\everypar{}

\parshape=0
\markboth{1 John 3:18}{1 John 3:18}\vs{18} Little children, let us not love with word or with tongue but in deed and truth.
\markboth{1 John 3:19}{1 John 3:19}\vs{19} And by this we will know that we are of the truth and will convince our conscience in his presence,
\markboth{1 John 3:20}{1 John 3:20}\vs{20} that if our conscience condemns us, that God is greater than our conscience and knows all things.\gva{r}\marginnote{\gva{r}\,If an evil conscience convicts us, much more ought the judgment of God condemn us, who knows our hearts better than we ourselves do.}
\markboth{1 John 3:21}{1 John 3:21}\vs{21} Dear friends, if our conscience does not condemn us, we have confidence in the presence of God,
\markboth{1 John 3:22}{1 John 3:22}\vs{22} and whatever we ask we receive from him, because we keep his commandments and do the things that are pleasing to him.
\markboth{1 John 3:23}{1 John 3:23}\vs{23} Now this is his commandment: that we believe in the name of his Son Jesus Christ and love one another, just as he gave us the commandment.
\markboth{1 John 3:24}{1 John 3:24}\vs{24} And the person who keeps his commandments resides in God, and God in him. Now by this we know that God resides in us: by the Spirit he has given us.\gva{f}\marginnote{\gva{f}\,He means the Spirit of sanctification, whereby we are born again and live to God.}

\markboth{1 John 4:1}{1 John 4:1}
\Needspace*{8\baselineskip}
\ch{4} \allowchapbreak\hypertarget{ch-1john-4}{}\lettrine[lines=5]{\color{general}D}{\textsc{ear}} friends, do not believe every spirit, but test the spirits to determine if they are from God, because many false prophets have gone out into the world.\gva{a}\marginnote{\gva{a}\,This is spoken by metonymy and it is as if he had said, Believe not every one who says that he has a gift of the Holy Spirit to do the office of a prophet.}\everypar{}
\markboth{1 John 4:2}{1 John 4:2}\vs{2} By this you know the Spirit of God: Every spirit that confesses Jesus as the Christ who has come in the flesh is from God,\gva{b}\marginnote{\gva{b}\,He speaks simply of the doctrine, and not of the person.}\gva{c}\marginnote{\gva{c}\,The true Messiah.}\gva{d}\marginnote{\gva{d}\,Is true man.}
\markboth{1 John 4:3}{1 John 4:3}\vs{3} but every spirit that refuses to confess Jesus, that spirit is not from God, and this is the spirit of the antichrist, which you have heard is coming, and now is already in the world.
\everypar{}

\parshape=0
\markboth{1 John 4:4}{1 John 4:4}\vs{4} You are from God, little children, and have conquered them because the one who is in you is greater than the one who is in the world.
\markboth{1 John 4:5}{1 John 4:5}\vs{5} They are from the world; therefore they speak from the world's perspective and the world listens to them.
\markboth{1 John 4:6}{1 John 4:6}\vs{6} We are from God; the person who knows God listens to us, but whoever is not from God does not listen to us. By this we know the Spirit of truth and the spirit of deceit.\gva{e}\marginnote{\gva{e}\,True prophets, against whom are false prophets, that is, those who err and lead others into error.}
\everypar{}

\parshape=0
\markboth{1 John 4:7}{1 John 4:7}\vs{7} Dear friends, let us love one another, because love is from God, and everyone who loves has been fathered by God and knows God.
\markboth{1 John 4:8}{1 John 4:8}\vs{8} The person who does not love does not know God because God is love.
\markboth{1 John 4:9}{1 John 4:9}\vs{9} By this the love of God is revealed in us: that God has sent his one and only Son into the world so that we may live through him.
\markboth{1 John 4:10}{1 John 4:10}\vs{10} In this is love: not that we have loved God, but that he loved us and sent his Son to be the atoning sacrifice for our sins.
\everypar{}

\parshape=0
\markboth{1 John 4:11}{1 John 4:11}\vs{11} Dear friends, if God so loved us, then we also ought to love one another.
\markboth{1 John 4:12}{1 John 4:12}\vs{12} No one has seen God at any time. If we love one another, God resides in us, and his love is perfected in us.\gva{g}\marginnote{\gva{g}\,Is surely in us indeed, and in truth.}
\markboth{1 John 4:13}{1 John 4:13}\vs{13} By this we know that we reside in God and he in us: in that he has given us of his Spirit.
\markboth{1 John 4:14}{1 John 4:14}\vs{14} And we have seen and testify that the Father has sent the Son to be the Savior of the world.
\everypar{}

\parshape=0
\markboth{1 John 4:15}{1 John 4:15}\vs{15} If anyone confesses that Jesus is the Son of God, God resides in him and he in God.\gva{h}\marginnote{\gva{h}\,With such a confession as comes from true faith, and is accompanied with love, so that there is an agreement of all things.}
\markboth{1 John 4:16}{1 John 4:16}\vs{16} And we have come to know and to believe the love that God has in us. God is love, and the one who resides in love resides in God, and God resides in him.
\markboth{1 John 4:17}{1 John 4:17}\vs{17} By this love is perfected with us, so that we may have confidence in the day of judgment, because just as Jesus is, so also are we in this world.\gva{i}\marginnote{\gva{i}\,This signifies a likeness, not an equality.}
\markboth{1 John 4:18}{1 John 4:18}\vs{18} There is no fear in love, but perfect love drives out fear because fear has to do with punishment. The one who fears punishment has not been perfected in love.\gva{k}\marginnote{\gva{k}\,If we understand by love, that we are in God, and God in us, that we are sons, and that we know God, and that everlasting life is in us: he concludes correctly, that we may well gather peace and quietness by this.}
\markboth{1 John 4:19}{1 John 4:19}\vs{19} We love because he loved us first.
\everypar{}

\parshape=0
\markboth{1 John 4:20}{1 John 4:20}\vs{20} If anyone says ``I love God'' and yet hates his fellow Christian, he is a liar because the one who does not love his fellow Christian whom he has seen cannot love God whom he has not seen.
\markboth{1 John 4:21}{1 John 4:21}\vs{21} And the commandment we have from him is this: that the one who loves God should love his fellow Christian too.

\markboth{1 John 5:1}{1 John 5:1}
\Needspace*{8\baselineskip}
\ch{5} \allowchapbreak\hypertarget{ch-1john-5}{}\lettrine[lines=5]{\color{general}E}{\textsc{veryone}} who believes that Jesus is the Christ has been fathered by God, and everyone who loves the father loves the child fathered by him.\gva{a}\marginnote{\gva{a}\,Is the true Messiah.}\gva{b}\marginnote{\gva{b}\,By one, he means all the faithful.}\everypar{}
\markboth{1 John 5:2}{1 John 5:2}\vs{2} By this we know that we love the children of God: whenever we love God and obey his commandments.\gva{c}\marginnote{\gva{c}\,There is no love where there is no true doctrine.}
\markboth{1 John 5:3}{1 John 5:3}\vs{3} For this is the love of God: that we keep his commandments. And his commandments do not weigh us down,\gva{d}\marginnote{\gva{d}\,To those who are regenerate, that is to say, born again, who are led by the Spirit of God, and are through grace delivered from the curse of the law.}
\everypar{}

\parshape=0
\markboth{1 John 5:4}{1 John 5:4}\vs{4} because everyone who has been fathered by God conquers the world. This is the conquering power that has conquered the world: our faith.\gva{e}\marginnote{\gva{e}\,He uses the time that is past, to give us to understand, that although we are in the battle, yet undoubtedly we shall be conquerors, and are most certain of the victory.}\gva{f}\marginnote{\gva{f}\,Which is the instrumental cause, and as a means and hand by which we lay hold on him, who indeed performs this, that is, has and does overcome the world, even Christ Jesus.}
\markboth{1 John 5:5}{1 John 5:5}\vs{5} Now who is the person who has conquered the world except the one who believes that Jesus is the Son of God?
\markboth{1 John 5:6}{1 John 5:6}\vs{6} Jesus Christ is the one who came by water and blood---not by the water only, but by the water and the blood. And the Spirit is the one who testifies, because the Spirit is the truth.\gva{g}\marginnote{\gva{g}\,Our spirit which is the third witness, testifies that the holy Sprit is truth, that is to say, that that is true which he tells us, that is, that we are the sons of God.}
\markboth{1 John 5:7}{1 John 5:7}\vs{7} For there are three that testify,\gva{h}\marginnote{\gva{h}\,See John 8:13-14}\gva{i}\marginnote{\gva{i}\,Agree in one.}
\markboth{1 John 5:8}{1 John 5:8}\vs{8} the Spirit and the water and the blood, and these three are in agreement.
\everypar{}

\parshape=0
\markboth{1 John 5:9}{1 John 5:9}\vs{9} If we accept the testimony of men, the testimony of God is greater because this is the testimony of God that he has testified concerning his Son.\gva{k}\marginnote{\gva{k}\,I conclude correctly: for the testimony which I said is given in heaven, comes from God, who sets forth his Son.}
\markboth{1 John 5:10}{1 John 5:10}\vs{10} (The one who believes in the Son of God has the testimony in himself; the one who does not believe God has made him a liar because he has not believed in the testimony that God has testified concerning his Son.)
\markboth{1 John 5:11}{1 John 5:11}\vs{11} And this is the testimony: God has given us eternal life, and this life is in his Son.
\markboth{1 John 5:12}{1 John 5:12}\vs{12} The one who has the Son has this eternal life; the one who does not have the Son of God does not have this eternal life.
\everypar{}

\parshape=0
\markboth{1 John 5:13}{1 John 5:13}\vs{13} I have written these things to you who believe in the name of the Son of God so that you may know that you have eternal life.
\everypar{}

\parshape=0
\markboth{1 John 5:14}{1 John 5:14}\vs{14} And this is the confidence that we have before him: that whenever we ask anything according to his will, he hears us.
\markboth{1 John 5:15}{1 John 5:15}\vs{15} And if we know that he hears us in regard to whatever we ask, then we know that we have the requests that we have asked from him.
\markboth{1 John 5:16}{1 John 5:16}\vs{16} If anyone sees his fellow Christian committing a sin not resulting in death, he should ask, and God will grant life to the person who commits a sin not resulting in death. There is a sin resulting in death. I do not say that he should ask about that.\gva{l}\marginnote{\gva{l}\,This is as if he said, let him ask the Lord to forgive him, and he will forgive him being so asked.}
\markboth{1 John 5:17}{1 John 5:17}\vs{17} All unrighteousness is sin, but there is sin not resulting in death.
\everypar{}

\parshape=0
\markboth{1 John 5:18}{1 John 5:18}\vs{18} We know that everyone fathered by God does not sin, but God protects the one he has fathered, and the evil one cannot touch him.
\markboth{1 John 5:19}{1 John 5:19}\vs{19} We know that we are from God, and the whole world lies in the power of the evil one.
\markboth{1 John 5:20}{1 John 5:20}\vs{20} And we know that the Son of God has come and has given us insight to know him who is true, and we are in him who is true, in his Son Jesus Christ. This one is the true God and eternal life.\gva{m}\marginnote{\gva{m}\,The divinity of Christ is most clearly proved by this passage.}
\markboth{1 John 5:21}{1 John 5:21}\vs{21} Little children, guard yourselves from idols.

\end{scripture}
